% Long \code{} run names are unbreakable; \sloppy keeps them from producing
% overfull boxes.  Scoped to this section only.
\begingroup
\sloppy

\section{Filter smoothing (hybrid)}
\label{sec:fs}

\newcommand{\HybMatrixTable}{%
\footnotesize
\begin{adjustbox}{max width=\textwidth}
\begin{tabular}{@{}l l l c c l l@{}}
\toprule
ID & truth & case & loc & $\Delta t_{\mathrm{E}}$ (s) & status & run directory \\
\midrule
H1 & PALM & inflow & on & 15 & \bad{failed (truth run)} & \code{pypalm\_to\_pyudales\_w3\_loccorrelation\_obs15\_inflow} \\
H2 & PALM & inflow\_turb & on & 15 & ok & \code{pypalm\_to\_pyudales\_w3\_loccorrelation\_obs15\_inflow\_turb} \\
H3 & PALM & periodic & on & 15 & ok & \code{pypalm\_to\_pyudales\_w3\_loccorrelation\_obs15\_periodic} \\
H4 & PALM & inflow\_turb & on & 30 & ok & \code{pypalm\_to\_pyudales\_w3\_loccorrelation\_obs30\_inflow\_turb} \\
H5 & PALM & periodic & on & 30 & ok & \code{pypalm\_to\_pyudales\_w3\_loccorrelation\_obs30\_periodic} \\
H6 & PALM & inflow\_turb & on & 7.5 & ok & \code{pypalm\_to\_pyudales\_w3\_loccorrelation\_obs7p5\_inflow\_turb} \\
H7 & PALM & periodic & on & 7.5 & ok & \code{pypalm\_to\_pyudales\_w3\_loccorrelation\_obs7p5\_periodic} \\
H8 & PALM & inflow & off & 15 & \bad{failed (truth run)} & \code{pypalm\_to\_pyudales\_w3\_locnone\_obs15\_inflow} \\
H9 & PALM & inflow\_turb & off & 15 & ok & \code{pypalm\_to\_pyudales\_w3\_locnone\_obs15\_inflow\_turb} \\
H10 & PALM & periodic & off & 15 & ok & \code{pypalm\_to\_pyudales\_w3\_locnone\_obs15\_periodic} \\
H11 & uDALES & inflow & on & 15 & ok & \code{pyudales\_to\_pyudales\_w3\_loccorrelation\_obs15\_inflow} \\
H12 & uDALES & inflow\_turb & on & 15 & ok & \code{pyudales\_to\_pyudales\_w3\_loccorrelation\_obs15\_inflow\_turb} \\
H13 & uDALES & periodic & on & 15 & ok & \code{pyudales\_to\_pyudales\_w3\_loccorrelation\_obs15\_periodic} \\
H14 & uDALES & inflow\_turb & on & 30 & ok & \code{pyudales\_to\_pyudales\_w3\_loccorrelation\_obs30\_inflow\_turb} \\
H15 & uDALES & periodic & on & 30 & ok & \code{pyudales\_to\_pyudales\_w3\_loccorrelation\_obs30\_periodic} \\
H16 & uDALES & inflow\_turb & on & 7.5 & ok & \code{pyudales\_to\_pyudales\_w3\_loccorrelation\_obs7p5\_inflow\_turb} \\
H17 & uDALES & periodic & on & 7.5 & ok & \code{pyudales\_to\_pyudales\_w3\_loccorrelation\_obs7p5\_periodic} \\
H18 & uDALES & inflow & off & 15 & ok & \code{pyudales\_to\_pyudales\_w3\_locnone\_obs15\_inflow} \\
H19 & uDALES & inflow\_turb & off & 15 & ok & \code{pyudales\_to\_pyudales\_w3\_locnone\_obs15\_inflow\_turb} \\
H20 & uDALES & periodic & off & 15 & ok & \code{pyudales\_to\_pyudales\_w3\_locnone\_obs15\_periodic} \\
\bottomrule
\end{tabular}
\end{adjustbox}
}

% AUTO-GENERATED by experiments_report/scripts/extract_filter_smoothing.py
\newcommand{\HybMasterTable}{%
\footnotesize
\begin{adjustbox}{max width=\textwidth}
\begin{tabular}{@{}l l c c r r r r r r r r r@{}}
\toprule
ID & case & loc & $\Delta t_{\mathrm{E}}$ & assim & valid & field $\Umag$ & $\alpha$ RMSE & $\Umag$ RMSE & $\chi^2$ & $z_{\mathrm{par}}$ & $q$ & wall \\
 & & & (s) & RMSE & RMSE & RMSE & [deg] (red \%) & [m/s] (red \%) & & & & [h] \\
\midrule
\multicolumn{13}{@{}l}{\textit{uDALES truth}}\\
H11 & inflow & on & 15 & 0.506 & 0.593 & 0.486 & 5.81 (66) & 0.616 (24) & 13.3 & 10.6 & 0.433 & 2.87 \\
H18 & inflow & off & 15 & 0.524 & 0.620 & 0.438 & 7.87 (53) & 0.506 (38) & 15.0 & 15.7 & 0.433 & 2.81 \\
H16 & inflow\_turb & on & 7.5 & 0.262 & 0.471 & 0.539 & \textbf{1.94} (88) & 0.500 (42) & 3.27 & 7.07 & 0.498 & 2.91 \\
H12 & inflow\_turb & on & 15 & 0.272 & 0.497 & 0.533 & 2.61 (84) & 0.433 (46) & 3.68 & 5.01 & 0.535 & 2.93 \\
H14 & inflow\_turb & on & 30 & 0.271 & 0.487 & 0.548 & 2.86 (83) & 0.512 (37) & 3.68 & 5.24 & 0.628 & 2.91 \\
H19 & inflow\_turb & off & 15 & \textbf{0.256} & \textbf{0.441} & \textbf{0.430} & 3.45 (82) & \textbf{0.365} (54) & 3.43 & 9.06 & 0.728 & 2.89 \\
H17 & periodic & on & 7.5 & 0.309 & 1.229 & 1.018 & 13.43 (18) & 0.738 (10) & 1.80 & 3.83 & 0.786 & 2.34 \\
H13 & periodic & on & 15 & 0.314 & 1.246 & 1.019 & 14.58 (12) & 0.756 (7) & 1.81 & 2.67 & 0.867 & 2.38 \\
H15 & periodic & on & 30 & 0.312 & 1.271 & 1.023 & 14.69 (12) & 0.764 (7) & 1.91 & \textbf{2.59} & 0.709 & 2.30 \\
H20 & periodic & off & 15 & 0.706 & 1.766 & 1.529 & 8.89 (45) & 0.928 (-27) & 23.1 & 41.5 & 0.761 & 2.35 \\
\midrule
\multicolumn{13}{@{}l}{\textit{PALM truth (model error)}}\\
H1 & inflow & on & 15 & \multicolumn{9}{c}{\bad{failed (truth run)}} \\
H8 & inflow & off & 15 & \multicolumn{9}{c}{\bad{failed (truth run)}} \\
H6 & inflow\_turb & on & 7.5 & 0.852 & 1.473 & \textbf{0.987} & 8.63 (42) & \textbf{0.414} (45) & 37.9 & 15.5 & -- & 2.85 \\
H2 & inflow\_turb & on & 15 & 0.812 & 1.429 & 0.991 & 6.96 (57) & 0.435 (44) & 34.4 & 12.4 & -- & 2.93 \\
H4 & inflow\_turb & on & 30 & 0.816 & 1.445 & 1.002 & 8.89 (39) & 0.485 (36) & 33.3 & 10.5 & -- & 2.89 \\
H9 & inflow\_turb & off & 15 & 0.942 & \textbf{1.416} & 1.280 & \textbf{6.89} (56) & 0.458 (39) & 43.3 & 15.4 & -- & 2.96 \\
H7 & periodic & on & 7.5 & 0.623 & 1.696 & 4.262 & 14.45 (13) & 1.310 (-28) & 9.64 & 4.86 & -- & 2.68 \\
H3 & periodic & on & 15 & \textbf{0.568} & 1.751 & 4.164 & 15.69 (6) & 1.157 (-19) & 4.35 & 2.97 & -- & 2.78 \\
H5 & periodic & on & 30 & 0.581 & 1.882 & 4.273 & 18.23 (2) & 1.386 (-27) & 3.86 & \textbf{2.86} & -- & 2.75 \\
H10 & periodic & off & 15 & 1.127 & 2.632 & 4.621 & 18.45 (-15) & 1.551 (-35) & 49.0 & 102 & -- & 2.80 \\
\bottomrule
\end{tabular}
\end{adjustbox}
}

\newcommand{\HybCalibTable}{%
\footnotesize
\begin{adjustbox}{max width=\textwidth}
\begin{tabular}{@{}l l c c r r r r r r r r r r r@{}}
\toprule
ID & case & loc & $\Delta t_{\mathrm{E}}$ & $\chi^2$ & $z_\alpha$ & $z_{\Umag}$ & $z_{\mathrm{pool}}$ & $c_\alpha$ & $c_{\Umag}$ & $q$ & $q_u$ & $q_v$ & $q_w$ & edge \\
 & & & (s) & & & & & & & & & & & frac \\
\midrule
\multicolumn{15}{@{}l}{\textit{uDALES truth}}\\
H11 & inflow & on & 15 & 13.3 & 11.3 & 10.0 & 10.6 & 0.32 & 0.35 & 0.433 & 0.921 & 0.161 & 0.217 & 0.684 \\
H18 & inflow & off & 15 & 15.0 & 17.6 & 14.0 & 15.7 & 0.33 & 0.33 & 0.433 & 0.925 & 0.143 & 0.230 & 0.730 \\
H16 & inflow\_turb & on & 7.5 & 3.27 & 4.65 & 8.53 & 7.07 & 0.24 & 0.33 & 0.498 & 0.932 & 0.300 & 0.263 & 0.280 \\
H12 & inflow\_turb & on & 15 & 3.68 & 3.77 & 6.12 & 5.01 & 0.25 & 0.36 & 0.535 & 0.931 & 0.398 & 0.276 & 0.320 \\
H14 & inflow\_turb & on & 30 & 3.68 & 3.05 & 6.71 & 5.24 & 0.27 & 0.36 & 0.628 & 0.927 & 0.683 & 0.273 & 0.297 \\
H19 & inflow\_turb & off & 15 & 3.43 & 9.39 & 8.87 & 9.06 & 0.23 & 0.28 & 0.728 & 0.942 & 0.722 & 0.522 & 0.342 \\
H17 & periodic & on & 7.5 & 1.80 & 2.99 & 4.54 & 3.83 & 0.63 & 0.62 & 0.786 & 0.939 & 0.717 & 0.701 & 0.091 \\
H13 & periodic & on & 15 & 1.81 & 2.72 & 2.72 & 2.67 & 0.70 & 0.76 & 0.867 & 0.952 & 0.917 & 0.731 & 0.095 \\
H15 & periodic & on & 30 & 1.91 & 2.58 & 2.61 & 2.59 & 0.75 & 0.79 & 0.709 & 0.937 & 0.488 & 0.701 & 0.102 \\
H20 & periodic & off & 15 & 23.1 & 37.0 & 40.3 & 41.5 & 0.23 & 0.22 & 0.761 & 0.925 & 0.910 & 0.449 & 0.625 \\
\midrule
\multicolumn{15}{@{}l}{\textit{PALM truth (model error)}}\\
H6 & inflow\_turb & on & 7.5 & 37.9 & 20.5 & 7.14 & 15.5 & 0.23 & 0.35 & -- & -- & -- & -- & 0.705 \\
H2 & inflow\_turb & on & 15 & 34.4 & 15.9 & 6.49 & 12.4 & 0.25 & 0.35 & -- & -- & -- & -- & 0.672 \\
H4 & inflow\_turb & on & 30 & 33.3 & 13.3 & 6.28 & 10.5 & 0.25 & 0.38 & -- & -- & -- & -- & 0.703 \\
H9 & inflow\_turb & off & 15 & 43.3 & 17.5 & 13.0 & 15.4 & 0.23 & 0.28 & -- & -- & -- & -- & 0.799 \\
H7 & periodic & on & 7.5 & 9.64 & 3.19 & 5.40 & 4.86 & 0.61 & 0.72 & -- & -- & -- & -- & 0.239 \\
H3 & periodic & on & 15 & 4.35 & 3.35 & 2.62 & 2.97 & 0.69 & 0.77 & -- & -- & -- & -- & 0.213 \\
H5 & periodic & on & 30 & 3.86 & 3.18 & 2.62 & 2.86 & 0.63 & 0.85 & -- & -- & -- & -- & 0.207 \\
H10 & periodic & off & 15 & 49.0 & 92.1 & 92.5 & 102 & 0.21 & 0.21 & -- & -- & -- & -- & 0.764 \\
\bottomrule
\end{tabular}
\end{adjustbox}
}

\newcommand{\HybNumOk}{18}
\newcommand{\HybNumRuns}{20}
\newcommand{\HybWallMin}{2.30}
\newcommand{\HybWallMax}{2.96}
\newcommand{\HybCycleMin}{46}
\newcommand{\HybCycleMax}{59}
\newcommand{\HybZMin}{2.59}
\newcommand{\HybZMax}{102}


\subsection{Method and experiment matrix}
\label{sec:fs:matrix}

Filter smoothing splits the two estimation problems between two estimators
inside every assimilation window. \textbf{Parameters} are updated by
\code{TimeVaryingParameterESMDA} (three MDA steps, \code{num\_time\_points=1},
\code{pin\_initial\_time\_point=true}) run against \emph{aggregated}
observations: the mean over bins of \code{esmda.interval\_seconds}.
\textbf{State} is updated by a stochastic \code{EnsembleKalmanFilter}
(\code{filtering.mode=state}, \code{StochasticEnKFAnalysis}, RTPS inflation
$\alpha = 0.6$) that assimilates \emph{every} $2$~s cycle. The smoother is
called with \code{final\_forecast=False}, so the window posterior state comes
from the filter, not from the last MDA iteration. Everything else --- domain,
ensemble size $N = 50$, $\sigma_{\mathrm{obs}} = 0.1$, three $120$~s windows to
$t = 360$~s, six assimilated and four validation sensors, $15$ parameter knots,
the AR(2) climatological prior --- is the shared setup of Section~\ref{sec:setup}.

Two knobs are swept on top of the three boundary-condition cases: localization
on/off (\code{loccorrelation} with truncation $0.35$, tapering $\beta = 0.5$,
max inflation $8$, applied to \emph{both} the filter and the smoother, versus
\code{locnone}), and the \textbf{ESMDA aggregation interval}
$\Delta t_{\mathrm{E}} = $ \code{esmda.interval\_seconds} $\in \{7.5, 15,
30\}$~s.

\begin{quote}
\textbf{\bad{Caveat on the method knob.}} $\Delta t_{\mathrm{E}}$ changes
\emph{only} the width of the bin the smoother averages observations over ($192$,
$96$ and $48$ aggregated observations per window). \textbf{The EnKF still
assimilates every $2$~s cycle in all \HybNumRuns{} runs.} This is not an
observation-frequency study and must not be read as one
(\S\ref{sec:fs:interval}).
\end{quote}

\begin{table}[H]
  \centering
  \caption{Filter smoothing: the \HybNumRuns{} experiments, their run IDs and
  the run directories under
  \code{presentations/isda\_new/} \code{experiments/filter\_smoothing/}. IDs are
  assigned by sorting the directory names alphabetically; the directory name is
  given here and nowhere else in this section. \code{loc} is
  \code{loccorrelation} (on) or \code{locnone} (off). Reference run of each
  case: H11, H12, H13.}
  \label{tab:fs:matrix}
  \HybMatrixTable
\end{table}

\subsection{Master results}
\label{sec:fs:master}

\begin{table}[H]
  \centering
  \caption{Filter smoothing: all \HybNumRuns{} runs, uDALES-truth block first.
  ``assim''/``valid'' are the vector RMSEs at the six assimilated and four
  validation sensors; ``field $\Umag$'' is the volume velocity-magnitude RMSE;
  $\alpha$ and $\Umag$ are parameter RMSEs with the reduction vs prior in
  parentheses; $\chi^2$ is the mean innovation chi-square over the $180$ cycles
  ($1$ = calibrated); $z_{\mathrm{par}}$ is the pooled parameter $z$-score
  standard deviation ($1$ = calibrated, $\gg 1$ = overconfident); $q$ is the
  field hit rate, undefined for PALM truth because truth and model live on
  different grids. Bold marks the best value per column within each truth
  block.}
  \label{tab:fs:master}
  \HybMasterTable
\end{table}

Read Table~\ref{tab:fs:master} case by case, not row by row: the boundary
condition, not the tuning, decides what works. State and observation
performance is best on \code{periodic} (assimilated-sensor vector RMSE
$0.31$, field hit rate up to $0.867$, $\chi^2 \approx 1.8$); parameter recovery
is best on \code{inflow\_turb} ($\alpha$ RMSE down to $1.94^\circ$, an $88\%$
reduction vs prior); and no run is good at both. The laminar \code{inflow} case,
which reads as the easy one, is the worst calibrated of the three. The
$\Delta t_{\mathrm{E}}$ column moves almost nothing --- as it must, given the
caveat above. Localization off buys a lower mean $\alpha$ error on
\code{periodic} while wrecking every other column, which is the signature of a
collapsed ensemble rather than a better estimate. Under PALM truth the sensor
fit degrades by roughly a factor of $2$--$3$ and both parameter reductions turn
negative on \code{periodic}.

\subsection{Results by case}
\label{sec:fs:cases}

All three subsections use the reference run of the case: uDALES truth,
localization on, $\Delta t_{\mathrm{E}} = 15$~s.

\subsubsection{Laminar inflow (\code{inflow})}
\label{sec:fs:inflow}

\begin{figure}[H]
  \centering
  \includegraphics[width=\textwidth]{filter_smoothing/valid_H11.png}
  \caption{H11 (uDALES truth, \code{inflow}, localization on,
  $\Delta t_{\mathrm{E}} = 15$~s): validation sensor~0 ($z = 2$~m), validation
  sensor~3 ($z = 20$~m) and the all-validation-sensor error panel. The x-axis is
  the \emph{cycle index} ($0$--$180$, $2$~s per cycle), not seconds as in the
  ESMDA and filtering sensor figures. The 50 members are visually
  indistinguishable from their mean --- the ensemble has collapsed --- and the
  mean sits below the truth for long stretches, a bias rather than noise.}
  \label{fig:fs:valid_H11}
\end{figure}

\begin{figure}[H]
  \centering
  \includegraphics[width=0.92\textwidth]{filter_smoothing/parerr_H11.png}
  \caption{H11: per-knot parameter error, inflow angle $\alpha$ above and
  velocity magnitude $\Umag$ below (blue RMSE, orange CRPS, shaded band =
  window~2). This figure is in seconds. Window~1 is converged to below
  $1.5^\circ$; the error then grows monotonically to $15.8^\circ$ at
  $t = 300$~s, and RMSE and CRPS nearly coincide throughout --- an ensemble with
  no spread left to score.}
  \label{fig:fs:parerr_H11}
\end{figure}

\begin{figure}[H]
  \centering
  \includegraphics[width=0.82\textwidth]{filter_smoothing/marginals_H11.png}
  \caption{H11: prior and posterior marginals at the final parameter knot (14 of
  15), $\alpha$ left and $\Umag$ right; dashed line = truth ($18.8^\circ$,
  $8.08$~m/s). The posterior overshoots to $28^\circ$ and $9.1$~m/s with a
  spread far too narrow to cover the truth ($z = -2.60$ and $-4.80$).}
  \label{fig:fs:marginals_H11}
\end{figure}

\begin{itemize}[nosep]
  \item \textbf{State.} Assimilated-sensor vector RMSE $0.506$, validation
        $0.593$, field $\Umag$ RMSE $0.486$ (H11; $0.524$/$0.620$/$0.438$
        without localization, H18). This is the \emph{worst} non-periodic
        validation result in the uDALES block --- the hybrid does not win the
        easy case.
  \item \textbf{Parameters.} $\alpha$ RMSE $5.81^\circ$ ($66\%$ reduction vs
        prior) and $\Umag$ RMSE $0.616$~m/s ($24\%$). The run-mean looks
        respectable only because window~1 is nearly perfect; by the final knot
        the angle error is $9.57^\circ$ (Fig.~\ref{fig:fs:parerr_H11}).
  \item \textbf{Calibration.} \bad{The worst of the three cases:} $\chi^2 =
        13.3$, pooled $z = 10.6$, field hit rate $q = 0.433$ (the lowest in the
        table), and contraction ratios $0.32$/$0.35$. The rank histogram
        (Fig.~\ref{fig:fs:rank}, left) puts $68\%$ of the validation ranks in the
        two end bins against a uniform expectation of $4\%$.
  \item \textbf{Mechanism.} The smooth laminar flow gives the filter almost
        nothing to correct, the ensemble contracts, and it then cannot track the
        parameter drift. Only $\Delta t_{\mathrm{E}} = 15$~s was run for this
        case, so there is no interval sweep here.
\end{itemize}

\subsubsection{Turbulent inflow (\code{inflow\_turb})}
\label{sec:fs:inflowturb}

\begin{figure}[H]
  \centering
  \includegraphics[width=\textwidth]{filter_smoothing/valid_H12.png}
  \caption{H12 (uDALES truth, \code{inflow\_turb}, localization on,
  $\Delta t_{\mathrm{E}} = 15$~s): validation sensors~0 and~3 plus the error
  panel; x-axis in cycle index ($2$~s per cycle). The ensemble keeps a visible
  spread and tracks the fluctuating $1.0$--$2.1$~m/s ground signal cycle by
  cycle; on-figure magnitude-only RMSE $0.367$ against a vector RMSE of $0.497$
  in Table~\ref{tab:fs:master}.}
  \label{fig:fs:valid_H12}
\end{figure}

\begin{figure}[H]
  \centering
  \includegraphics[width=0.92\textwidth]{filter_smoothing/parerr_H12.png}
  \caption{H12: per-knot $\alpha$ and $\Umag$ error (seconds on the x-axis).
  Unlike \code{inflow} the error stays bounded across all three windows and a
  healthy RMSE--CRPS gap survives to the end: this is the only case in which the
  parameters are genuinely recovered.}
  \label{fig:fs:parerr_H12}
\end{figure}

\begin{figure}[H]
  \centering
  \includegraphics[width=0.82\textwidth]{filter_smoothing/marginals_H12.png}
  \caption{H12: final-knot marginals. The posterior lands at $15.2^\circ$
  against a truth of $18.8^\circ$ --- much closer than either other case --- but
  the spread is so narrow that $z = 3.59$; $\Umag$ overshoots to $8.8$~m/s
  ($z = -6.28$). The final knot is the worst knot in every run.}
  \label{fig:fs:marginals_H12}
\end{figure}

\begin{itemize}[nosep]
  \item \textbf{Parameters --- the one case that works.} $\alpha$ RMSE
        \good{$1.94$--$3.45^\circ$} across the four matched-model runs
        ($82$--$88\%$ reduction vs prior), bounded in every window;
        $\Umag$ RMSE $0.365$--$0.512$~m/s ($37$--$54\%$).
  \item \textbf{State --- the best of the sweep.} H19 (localization off) has the
        lowest validation vector RMSE ($0.441$), the lowest assimilated-sensor
        RMSE ($0.256$) and the lowest field $\Umag$ RMSE ($0.430$) in the whole
        uDALES block.
  \item \textbf{Calibration --- mid-range and clearly imperfect.}
        $\chi^2 = 3.27$--$3.68$, pooled $z = 5.01$--$9.06$, $q = 0.498$--$0.728$,
        validation rank-edge fraction $0.28$--$0.34$. Under-dispersive, but
        nothing like \code{inflow}.
  \item \textbf{Localization is close to a wash here.} It helps the angle
        ($2.61^\circ$ on, $3.45^\circ$ off) and the pooled $z$ ($5.01$ vs
        $9.06$); it costs on the sensors ($0.497$ vs $0.441$) and on the hit
        rate ($0.535$ vs $0.728$).
\end{itemize}

\subsubsection{Periodic boundaries (\code{periodic})}
\label{sec:fs:periodic}

\begin{figure}[H]
  \centering
  \includegraphics[width=\textwidth]{filter_smoothing/valid_H13.png}
  \caption{H13 (uDALES truth, \code{periodic}, localization on,
  $\Delta t_{\mathrm{E}} = 15$~s): validation sensors~0 and~3 plus the error
  panel; x-axis in cycle index. This is the honest limitation of the best-state
  run: at sensors that were never assimilated the ensemble median is damped and
  the truth leaves the ensemble repeatedly (on-figure magnitude RMSE $0.688$),
  even though the assimilated sensors are fitted to $0.314$ vector RMSE.}
  \label{fig:fs:valid_H13}
\end{figure}

\begin{figure}[H]
  \centering
  \includegraphics[width=0.92\textwidth]{filter_smoothing/parerr_H13.png}
  \caption{H13: per-knot $\alpha$ and $\Umag$ error (seconds on the x-axis). The
  angle error never falls decisively below the $16.6^\circ$ climatological prior
  and ends at $25.8^\circ$ --- worse than the prior. The forcing is not
  identifiable from six near-ground sensors under periodic lateral boundaries.}
  \label{fig:fs:parerr_H13}
\end{figure}

\begin{figure}[H]
  \centering
  \includegraphics[width=0.82\textwidth]{filter_smoothing/marginals_H13.png}
  \caption{H13: final-knot marginals. The prior mean sits at $+2^\circ$ against a
  truth of $18.8^\circ$ and the update moves the posterior to $-6^\circ$, i.e.\
  \emph{away} from the truth ($z = 3.52$); $\Umag$ barely moves ($z = 1.32$).
  This is the figure that forbids presenting \code{periodic} as a
  parameter-estimation success.}
  \label{fig:fs:marginals_H13}
\end{figure}

\begin{itemize}[nosep]
  \item \textbf{State --- the best in the section.} Assimilated-sensor vector
        RMSE $0.309$--$0.314$ with localization; field $\Umag$ RMSE
        $1.018$--$1.023$. Validation is the weak side ($1.229$--$1.271$): the
        held-out sensors sit in a flow the six assimilated ones do not constrain.
  \item \textbf{Calibration --- also the best in the section.} $\chi^2 = 1.80$,
        $1.81$, $1.91$ for $\Delta t_{\mathrm{E}} = 7.5/15/30$~s; pooled $z$ down
        to $2.59$ (H15); $q = 0.867$ for H13 with $q_u = 0.952$, $q_v = 0.917$,
        $q_w = 0.731$; validation rank-edge fraction $0.09$--$0.10$, by far the
        closest to uniform anywhere in the sweep.
  \item \bad{\textbf{Parameters --- a failure.}} $\alpha$ RMSE $13.4$--$14.7^\circ$,
        i.e.\ a $12$--$18\%$ reduction against a $16.4$--$16.7^\circ$
        climatological prior, with a final-knot error of $25.8^\circ$ for H13;
        $\Umag$ improves by $7$--$10\%$.
  \item \textbf{And it is unidentifiability, not non-convergence.} The smoother
        drives its normalized data mismatch onto target on this case and the
        posterior observation RMSE reaches $0.053$ against
        $\sigma_{\mathrm{obs}} = 0.1$: it fits the aggregated observations almost
        perfectly with the \emph{wrong} boundary-condition parameters. The right
        framing for \code{periodic} is \emph{state estimation under an
        unidentifiable forcing}.
\end{itemize}

\subsection{Model error (PALM truth)}
\label{sec:fs:modelerror}

Rows H2--H10 of Table~\ref{tab:fs:master} replace the uDALES truth by PALM while
keeping uDALES as the assimilation model, so the ensemble carries genuine
structural model error.

\begin{figure}[H]
  \centering
  \includegraphics[width=0.95\textwidth]{filter_smoothing/modelerror_valid_inflow_turb.png}
  \caption{H12 (uDALES truth) versus H2 (PALM truth) on \code{inflow\_turb} at
  identical settings: validation sensor~0 and the validation-error panel, x-axis
  in cycle index. Under model error the ensemble is flat and stays near
  $1.3$~m/s while the truth swings between $0.5$ and $3.0$~m/s; the on-figure
  magnitude RMSE rises from $0.367$ to $0.974$.}
  \label{fig:fs:modelerror}
\end{figure}

\begin{itemize}[nosep]
  \item \textbf{\code{inflow\_turb} is the presentable half.} The angle still
        drops to $6.89$--$8.89^\circ$ ($39$--$57\%$ reduction vs prior) and
        $\Umag$ to $0.414$--$0.485$~m/s ($36$--$45\%$); field $\Umag$ RMSE
        ($0.987$--$1.28$) is about twice the matched-model value. But the
        calibration is the worst in the section: $\chi^2 = 33.3$--$43.3$, pooled
        $z = 10.5$--$15.5$, validation rank-edge fraction $0.67$--$0.80$.
  \item \bad{\textbf{\code{periodic} loses the field entirely.}} Field $\Umag$
        RMSE jumps to $4.16$--$4.27$ against $1.02$ for matched model --- a
        factor of $4$ --- and validation vector RMSE to $1.70$--$1.88$ from
        $1.25$, while the assimilated sensors are still fitted ($0.568$--$0.623$).
        The filter absorbs structural error into the state and the boundary
        condition instead of correcting it.
  \item \textbf{The parameter update reverses sign.} On every PALM
        \code{periodic} run $\Umag$ ends \emph{worse} than its prior ($-19$ to
        $-35\%$) and the angle reduction is $2$--$13\%$ (and $-15\%$ without
        localization) --- the classic model-error signature.
  \item \textbf{The worst run of the sweep is H10} (model error, \code{periodic},
        localization off): $\chi^2 = 49.0$, pooled $z = 102$, field $\Umag$ RMSE
        $4.62$, validation $2.63$, both parameters worse than prior.
  \item The field hit rate $q$ is undefined for the whole PALM block: truth and
        assimilation model are on different grids, so only sensor-space and
        parameter metrics are comparable across the two truth blocks.
\end{itemize}

\subsection{Sensitivities}
\label{sec:fs:sens}

\subsubsection{Localization}
\label{sec:fs:loc}

\begin{figure}[H]
  \centering
  \includegraphics[width=0.86\textwidth]{filter_smoothing/localization_parerr_periodic.png}
  \caption{Localization on/off on \code{periodic} at matched settings: H13 (top
  pair) versus H20 (bottom pair), $\alpha$ above $\Umag$ within each pair;
  seconds on the x-axis. In H20 the RMSE and CRPS curves lie on top of each other
  at every knot --- the textbook signature of a collapsed ensemble, whose lower
  mean angle error is an artefact of near-zero spread rather than a better
  estimate.}
  \label{fig:fs:loc}
\end{figure}

Five matched pairs differ only by localization (H11/H18, H12/H19, H13/H20,
H2/H9, H3/H10 in Table~\ref{tab:fs:master}). The effect is strongly
case-dependent, and on \code{periodic} it has opposite signs for the parameters
and for everything else.

\begin{itemize}[nosep]
  \item \code{inflow}: localization helps the angle ($5.81^\circ$ vs
        $7.87^\circ$) and the calibration (pooled $z$ $10.6$ vs $15.7$) and is
        neutral on the sensors ($0.593$ vs $0.620$ validation).
  \item \code{inflow\_turb}: a wash, as listed in \S\ref{sec:fs:inflowturb}.
  \item \code{periodic}: \bad{turning localization off reports a better mean
        angle error ($8.89^\circ$ vs $14.58^\circ$) that must not be believed.}
        The same run has $\chi^2 = 23.1$ (vs $1.81$), pooled $z = 41.5$ (vs
        $2.67$), $\Umag$ \emph{worse} than its prior ($-27\%$), validation vector
        RMSE $42\%$ worse ($1.766$ vs $1.246$), field $\Umag$ RMSE $50\%$ worse
        ($1.529$ vs $1.019$) and a validation rank-edge fraction of $0.625$ vs
        $0.095$.
  \item Under model error localization is unambiguously necessary: dropping it
        takes pooled $z$ from $2.97$ to $102$ on \code{periodic} and degrades
        every state metric.
\end{itemize}

\subsubsection{ESMDA aggregation interval}
\label{sec:fs:interval}

\begin{figure}[H]
  \centering
  \includegraphics[width=0.88\textwidth]{filter_smoothing/interval_parerr_inflow_turb.png}
  \caption{H16, H12 and H14 (uDALES truth, \code{inflow\_turb}, localization on,
  $\Delta t_{\mathrm{E}} = 7.5$, $15$, $30$~s): the inflow-angle error panel of
  each run, seconds on the x-axis. Note the differing vertical scales. Finer
  aggregation bins retain more information for the smoother --- the trend is
  real but only about $1^\circ$ wide, an order of magnitude smaller than the
  differences between cases.}
  \label{fig:fs:interval}
\end{figure}

Twelve runs form four $7.5/15/30$~s sweeps (two cases $\times$ two truth models).
Because only the smoother's aggregation bin changes --- the filter assimilates
every $2$~s in every row --- the observed sensitivity is exactly what that
implies.

\begin{itemize}[nosep]
  \item \textbf{State metrics are essentially inert.} \code{inflow\_turb} with
        uDALES truth: field $\Umag$ RMSE $0.539/0.533/0.548$, validation
        $0.471/0.497/0.487$. \code{periodic}: field $1.018/1.019/1.023$,
        validation $1.229/1.246/1.271$. Every spread is under $4\%$.
  \item \textbf{Parameters show a small monotone trend on \code{inflow\_turb}}:
        $\alpha$ RMSE $1.94/2.61/2.86^\circ$ (Fig.~\ref{fig:fs:interval}).
  \item \textbf{On \code{periodic} even the parameters barely move}:
        $13.4/14.6/14.7^\circ$, all within noise of the climatological prior.
  \item Under model error the trend is not even monotone (\code{inflow\_turb}:
        $8.63/6.96/8.89^\circ$; \code{periodic}: $14.5/15.7/18.2^\circ$),
        consistent with model error dominating the aggregation width.
\end{itemize}

\subsection{Calibration}
\label{sec:fs:calibration}

\begin{table}[H]
  \centering
  \caption{Filter smoothing: calibration diagnostics for the \HybNumOk{}
  completed runs. $z_\bullet$ are parameter $z$-score standard deviations
  ($1$ = calibrated); $c_\bullet$ are mean posterior/prior spread contraction
  ratios; $q$ is the field hit rate with its per-component values; ``edge frac''
  is the fraction of pooled validation-sensor ranks falling in the two extreme
  bins, against a uniform expectation of $0.039$ for $N = 50$.}
  \label{tab:fs:calib}
  \HybCalibTable
\end{table}

\begin{figure}[H]
  \centering
  \includegraphics[width=\textwidth]{filter_smoothing/rank_hist_cases.png}
  \caption{Posterior rank histograms of the truth within the 50-member ensemble
  for the three reference runs H11, H12 and H13 (top row: assimilated sensors,
  $N = 4320$; bottom row: validation sensors, $N = 2880$). Dashed line =
  uniform, grey band = $\pm 2\sigma$. \code{inflow} is violently U-shaped (end
  bins about $4\times$ uniform): gross under-dispersion. \code{inflow\_turb} is
  moderately U-shaped. \code{periodic} is near-uniform on the assimilated
  sensors with only a mild U on the validation sensors.}
  \label{fig:fs:rank}
\end{figure}

Three facts stand out of Table~\ref{tab:fs:calib}.

\begin{enumerate}[nosep]
  \item \textbf{Every run is parameter-overconfident.} The pooled $z$-score
        standard deviation ranges from \HybZMin{} to \HybZMax{} where $1$ is
        calibrated; even the best run (H15) is off by a factor of $2.6$.
        \bad{No ``well-calibrated parameters'' claim can be supported by these
        runs.}
  \item \textbf{The ordering ``\code{inflow} easy $\rightarrow$ \code{periodic}
        hard'' is false for calibration.} \code{inflow} has the worst $\chi^2$
        ($13.3$--$15.0$), the worst hit rate ($0.433$) and the worst rank-edge
        fraction ($0.68$--$0.73$); \code{periodic} with localization has the best
        of all three ($\chi^2 \approx 1.8$, $q$ up to $0.867$, edge fraction
        $0.09$). Difficulty for the state and difficulty for the parameters are
        anti-correlated here.
  \item \textbf{The vertical component is the weak one.} Where $q$ resolves by
        component, $q_u$ is $0.92$--$0.95$ everywhere, whereas $q_w$ ranges from
        $0.217$ (\code{inflow}) to $0.731$ (H13). $q_v$ is the most
        case-sensitive ($0.143$ to $0.917$).
\end{enumerate}

\subsection{Failed runs and cost}
\label{sec:fs:failed}

Two of the \HybNumRuns{} directories, H1 and H8 --- the only two PALM-truth
\code{inflow} runs --- contain only \code{config.yaml}, \code{hydra/} and an
empty \code{windows/}. Both failed for the same reason, recorded in
\code{\_logs/*.log}: PALM produced only $52$ of the $180$ expected $2$~s 3D
output frames, i.e.\ it diverged after roughly $104$~s of a $360$~s truth
trajectory, and \code{pypalm.forward\_model.\_fit\_output\_window} refused to pad
the missing $128$ frames. The failure happens while generating the truth
trajectory, before any assimilation, so this is a PALM stability problem in the
laminar-inflow configuration (disturbances and initial seeding both off), not a
data-assimilation failure; it is independent of the localization setting, and
the matched-model \code{inflow} runs and the PALM \code{inflow\_turb} and
\code{periodic} runs are unaffected.

Cost is flat: \HybWallMin--\HybWallMax~h per run across all \HybNumOk{}
completed runs (\HybCycleMin--\HybCycleMax~s per $2$~s cycle at $N = 50$ with
$10$ parallel processes), with \code{periodic} slightly cheaper than the two
inflow cases. Neither the case, the localization setting nor
$\Delta t_{\mathrm{E}}$ changes the cost measurably, and all $50$ members remain
distinct in every run.

\subsection{Takeaways}
\label{sec:fs:takeaways}

\begin{enumerate}[nosep]
  \item \textbf{Both halves, neither best.} On \code{periodic} at matched
        settings the hybrid H13 fits the assimilated sensors to $0.314$ vector
        RMSE against $1.091$ for ESMDA-only (E13), and keeps $\Umag$ parameter
        RMSE at $0.756$~m/s against $4.221$~m/s for the joint EnKF (F11). It does
        \emph{not} beat ESMDA-only on parameters ($14.58^\circ$ vs $13.62^\circ$)
        and is statistically indistinguishable from a state-only EnKF with RTPS
        on the state ($1.246$ vs $1.261$ validation; $1.019$ vs $1.029$ field
        $\Umag$); its advantage over that baseline is that it also returns a
        boundary-condition estimate. Section~\ref{sec:cmp} settles the ranking.
  \item \good{\textbf{Parameters work on \code{inflow\_turb} only.}} $\alpha$
        RMSE $1.94$--$3.45^\circ$ ($82$--$88\%$ reduction vs prior) across all
        four matched-model runs, bounded in every window, at a validation vector
        RMSE of $0.441$--$0.497$.
  \item \bad{\textbf{Periodic forcing is unidentifiable.}} $12$--$18\%$ reduction
        vs prior, a final-knot error of $25.8^\circ$ and a posterior that moves
        away from the truth --- even though the smoother fits the aggregated
        observations to $0.053$ against $\sigma_{\mathrm{obs}} = 0.1$.
  \item \bad{\textbf{The interval knob is a placebo}} for the state: it changes
        only the smoother's aggregation bin, and state metrics move by under
        $4\%$ across $7.5/15/30$~s. It buys about $1^\circ$ of angle accuracy on
        \code{inflow\_turb} and nothing measurable elsewhere.
  \item \textbf{Localization keeps the ensemble alive.} A better parameter number
        without it is a red flag: H20 reports $8.89^\circ$ vs H13's $14.58^\circ$
        from an ensemble with $\chi^2 = 23.1$ and pooled $z = 41.5$, and pays
        $42\%$ on validation sensors. Under model error, localization off pushes
        pooled $z$ to $102$.
  \item \textbf{Model error costs a factor of four.} PALM truth on
        \code{periodic} gives field $\Umag$ RMSE $4.16$--$4.27$ versus $1.02$
        matched, and both parameters end worse than their prior on every such
        run, while the observation fit stays at $\sigma_{\mathrm{obs}}$.
  \item \textbf{The easy case is the worst calibrated.} Every run is
        parameter-overconfident (pooled $z$ \HybZMin--\HybZMax), and laminar
        \code{inflow} --- not \code{periodic} --- has the worst $\chi^2$
        ($13.3$--$15.0$) and the worst hit rate ($0.433$). Cost is a flat
        \HybWallMin--\HybWallMax~h per run throughout.
\end{enumerate}

\clearpage
\subsection{Appendix: all runs}
\label{sec:fs:appendix}

\begin{figure}[H]
  \centering
  \includegraphics[height=0.80\textheight]{filter_smoothing/parerr_grid_all.png}
  \caption{Inflow-angle error per knot for all \HybNumOk{} completed runs
  (same order as Table~\ref{tab:fs:master}; blue RMSE, orange CRPS, shaded
  band = window~2, seconds on the x-axis). \textbf{Vertical scales differ between
  panels.} Two patterns read at a glance: the \code{inflow\_turb} panels stay
  flat on a scale that tops out below $8^\circ$ while every \code{periodic}
  panel needs $15$--$35^\circ$, and in the localization-off panels the RMSE and
  CRPS curves collapse onto each other.}
  \label{fig:fs:grid}
\end{figure}

\endgroup
